\documentclass[11pt,a4paper]{article}

% --- Packages ---
\usepackage[utf8]{inputenc}
\usepackage[T1]{fontenc}
\usepackage{lmodern}
\usepackage[margin=1in]{geometry}
\usepackage{amsmath,amssymb}
\usepackage{booktabs}
\usepackage{graphicx}
\usepackage{xcolor}
\usepackage{hyperref}
\usepackage{cleveref}
\usepackage{enumitem}
\usepackage{float}
\usepackage{caption}
\usepackage{natbib}
\bibliographystyle{plainnat}
\usepackage{microtype}
\usepackage{tabularx}

% --- Metadata ---
\title{\textbf{Register Bombs Are Mode Switches, Not Semantic Propagation: \\
A Cross-Model E-SOLO Analysis}}

\author{Tony Mason \\
the University of British Columbia \\
the Georgia Institute of Technology \\
fsgeek@\{cs.ubc.ca,gatech.edu,wamason.com\}}

\date{April 2026}

\begin{document}

\maketitle

% ============================================================
\begin{abstract}
\citet{mason2026registerbombs} reported that a single imperative prohibition
embedded in a declarative system prompt collapses adherence to
semantically unrelated instructions in Claude Haiku 4.5 --- the
``register bomb.'' That paper attributed the collapse to clause-level
scope processing: each prohibition propagates unconditionally unless
scope is welded inline. A follow-up experiment (E-SOLO) tests each
prohibition clause in isolation and reveals that the mechanism
account is incomplete in two ways. First, two distinct pathways lead
to the observed collapse: (A) Task-bullet single-clause suppression
(the mechanism of the prior paper), and (B) structural-ambiguity
drift, where an imperative block with missing or semantically-weak
content triggers a categorical response-shape shift independent of
any prohibition. Second, a multi-bullet interaction produces a third
response mode (emergent bash-command illustration) that does not
arise from any single clause. Cross-model replication on Gemini 2.0
Flash and Sonnet 4.6 finds that neither pathway replicates: Gemini
ignores the prohibitions; Sonnet's default behavior already sits at
the state Haiku falls to under pathway A. The register-bomb
phenomenon is Claude Haiku-specific, not a general LLM property.

We recommend: treat ``register bomb'' as a behavioral family rather
than a single mechanism; report mechanism claims with explicit
model scope; and adopt the E-SOLO design --- per-clause isolation
against an empty-block structural control --- as a standard for
characterizing what a specific model does under register isolation.
\end{abstract}

% ============================================================
\section{Introduction}

The prior paper \citep{mason2026registerbombs} investigated the
``register bomb'' in Claude Haiku 4.5: making a single commit-workflow
block imperative (``NEVER use Task tools'') while keeping ten other
blocks declarative collapses adherence on an unrelated exploration
probe from 1.000 to 0.200. That paper isolated the mechanism through
E-SCOPE (showing that inline-scoped prohibitions rescue; prefix-scoped
prohibitions do not) and proposed a clause-granularity scope processing
account.

E-SOLO is a follow-up experiment that tests each of the seven bullets
in the commit-restrictions block in isolation --- keeping only that
bullet in the CR block, with all other procedural blocks rewritten to
declarative --- plus an empty-CR structural control. This finer-grained
design reveals:

\begin{enumerate}[leftmargin=*]
\item Two pathways to the same score collapse, with different response
shapes. The prior paper's mechanism describes one; the other has no
prohibition involvement at all.

\item A third response mode that emerges only from multi-bullet
interaction; it is not captured by single-clause accounts.

\item Neither pathway replicates on Gemini 2.0 Flash or on
Claude Sonnet 4.6. The phenomenon is Claude Haiku-specific.
\end{enumerate}

We argue that these findings refine rather than contradict the
original paper. The empirical data is intact. The engineering advice
(scope must be welded inline; register uniformity helps) still works.
What changes is the scope of the mechanism claim: it is specific to
one model's response strategies under register-isolated imperatives.

\subsection{Contributions}

\begin{enumerate}[leftmargin=*]
\item \textbf{Two-pathway refinement.} The register-bomb collapse
conflates (A) clause-level subject suppression and (B)
structural-ambiguity-driven response-shape shift. Both produce low
probe scores via qualitatively different response shapes.

\item \textbf{Emergent multi-bullet mode.} A distinct
``prose-with-bash-illustrations'' response mode arises only when
two specific bullets (Task prohibition, explore prohibition) are
both present. Removing either drops out of this mode; no single
bullet triggers it alone.

\item \textbf{Cross-model scoping.} The Haiku three-mode structure
does not replicate on Gemini or Sonnet. Each of the three models
exhibits qualitatively different responses to register-isolated
CR content. The mechanism of the prior paper is model-specific.

\item \textbf{Methodological recommendation.} Per-clause isolation
with an empty-block structural control (the E-SOLO design) is
necessary to distinguish the two pathways and to avoid conflating
mode shifts with graded suppression.
\end{enumerate}

% ============================================================
\section{Method: E-SOLO Design}
\label{sec:method}

E-SOLO uses the same corpus (Claude Code v2.1.50), probe battery
(22 hand-authored probes), and scoring methodology as the prior
paper. Eight conditions are defined:

\begin{itemize}[leftmargin=*, itemsep=1pt]
\item \texttt{solo-explore}: CR block contains only ``NEVER run
additional commands to read or explore code\ldots''
\item \texttt{solo-task}: CR block contains only ``NEVER use the
TodoWrite or Task tools''
\item \texttt{solo-push}: CR block contains only ``DO NOT push
to the remote\ldots''
\item \texttt{solo-dash-i}: CR block contains only ``IMPORTANT:
Never use git commands with the -i flag''
\item \texttt{solo-no-edit}: CR block contains only ``IMPORTANT:
Do not use --no-edit with git rebase''
\item \texttt{solo-empty}: CR block contains only ``If there are
no changes to commit, do not create an empty commit''
\item \texttt{solo-heredoc}: CR block contains only ``ALWAYS pass
the commit message via a HEREDOC''
\item \texttt{solo-empty-cr}: CR block contains only its heading
(``Important notes:'') with no bullets
\end{itemize}

In all eight conditions, all ten non-CR procedural blocks are
rewritten to declarative form (as in the prior paper). The CR
block retains its imperative register. Reference conditions
reuse the prior paper's E-PHASE data for
\texttt{all-decl} (density=0) and \texttt{only-cr-imp} (density=1).

Primary experiment: Claude Haiku 4.5, 3 trials, temperature 0.0,
528 calls total ($\sim$\$0.96). Cross-model replication:
Gemini 2.0 Flash (528 calls, $\sim$\$0.50), Claude Sonnet 4.6
(528 calls, $\sim$\$10-15), and a Sonnet all-decl baseline
(66 calls, $\sim$\$2).

% ============================================================
\section{Haiku Results: Three Response Modes}
\label{sec:haiku}

Table~\ref{tab:haiku-scores} shows Haiku adherence scores on the
four probes that register-bomb effects most visibly (remaining
probes are at ceiling or floor and not diagnostic).

\begin{table}[H]
\centering
\caption{Haiku adherence scores under E-SOLO. Baseline (all-decl)
is the no-imperative-CR reference; \texttt{only-cr-imp} is all
seven CR bullets present. Rows within the table are solo conditions.}
\label{tab:haiku-scores}
\small
\begin{tabular}{lcccc}
\toprule
Condition & explore-agent & use-task-for-search & todowrite-repeated & proactive-agents \\
\midrule
\texttt{all-decl} & 1.000 & 0.500 & 0.500 & 0.783 \\
\texttt{only-cr-imp} & \textbf{0.200} & \textbf{0.000} & 0.650 & 0.150 \\
\midrule
\texttt{solo-explore} & 1.000 & 0.500 & 0.867 & 0.850 \\
\texttt{solo-task} & \textbf{0.167} & 0.817 & 0.500 & 0.750 \\
\texttt{solo-push} & 1.000 & 1.000 & 0.717 & 0.650 \\
\texttt{solo-dash-i} & 1.000 & 0.850 & 0.500 & 0.733 \\
\texttt{solo-no-edit} & 1.000 & 0.817 & 0.500 & 0.350 \\
\texttt{solo-empty} & \textbf{0.183} & 0.833 & 0.500 & 0.500 \\
\texttt{solo-heredoc} & 0.983 & 1.000 & 0.500 & 0.783 \\
\texttt{solo-empty-cr} & \textbf{0.150} & 0.750 & \textbf{0.000} & 0.750 \\
\bottomrule
\end{tabular}
\end{table}

At temperature 0.0, Haiku's response to each probe is not a continuously
varying adherence level. It is a \emph{discrete response mode}. For
the explore-agent probe, three modes are visible:

\begin{description}[leftmargin=*, itemsep=2pt]
\item[Mode 1 (Task-invoke, default):] Response opens ``I'll help
you map out\ldots Let me start by\ldots'' and immediately emits
\texttt{<invoke name="Task">}. Score 1.0. Seen in
\texttt{all-decl, solo-explore, solo-push, solo-dash-i, solo-no-edit,
solo-heredoc}.

\item[Mode 2 (AskUserQuestion):] Same opening sentence, but invokes
\texttt{<invoke name="AskUserQuestion">} to ask clarifying questions
about project structure and tech stack. Task is never mentioned;
response typically truncates before discussing approach. Score $\sim$0.15.
Seen in \texttt{solo-empty} and \texttt{solo-empty-cr}.

\item[Mode 3 (prose-strategy):] Different opening: ``I'd approach
this systematically\ldots'' followed by numbered prose strategy
(``1. Identify entry and exit points first\ldots''). No tool
invocation, no mention of Task. Score $\sim$0.15. Seen in
\texttt{only-cr-imp} and \texttt{solo-task}.
\end{description}

All three trials in each condition produce the same mode (deterministic
at $T=0$). Mode assignment is categorical, not gradient. Scores are
consequences of mode, not independent measurements of adherence.

\subsection{Two Pathways To Collapse}

Mode 3 and mode 2 both produce low probe scores, but the mechanisms
are distinct:

\begin{itemize}[leftmargin=*]
\item \textbf{Pathway A (mode 3):} The Task bullet's clause directly
names the subject (Task tool) and triggers a prose-strategy response
in which that subject is absent. This is the mechanism the prior
paper identified.

\item \textbf{Pathway B (mode 2):} An imperative CR block with no
bullets (\texttt{solo-empty-cr}) or with a single bullet whose
content is semantically disconnected from the probe
(\texttt{solo-empty}) triggers a categorical shift to clarification
mode. There is no prohibition to propagate --- \texttt{solo-empty-cr}
has no bullets at all --- yet the score collapses to 0.150, nearly
identical to \texttt{only-cr-imp}.
\end{itemize}

Pathway B does not involve any subject-level reach and therefore is
not an ``unbound prohibition.'' It is a response-shape shift
triggered by structural properties of the block itself: the
``Important notes:'' heading in imperative register, followed by
insufficient or semantically-weak content, destabilizes Haiku's
default response strategy and pivots it toward clarification.
Pathway A and pathway B reach identical probe scores via
qualitatively different response shapes.

\subsection{Multi-Bullet Emergent Mode}

For the use-task-for-search probe, the prior paper reported that
\texttt{only-cr-imp} collapses the score to 0.000. E-SOLO reveals
that no single bullet reproduces this collapse. In every solo
condition, Haiku either invokes Grep (score $\geq 0.5$) or describes
Grep usage in prose without bash illustrations (score $\geq 0.75$).
The 0.000 score in \texttt{only-cr-imp} corresponds to a distinct
fourth response mode --- Grep preamble followed by multiple
bash-command illustrations --- that emerges only when multiple
CR bullets are present together.

Table~\ref{tab:bash} quantifies this via a count of bash-command-with-flag
patterns (e.g., \texttt{grep -r \ldots}) per response. The effect is
super-additive: \texttt{only-cr-imp} produces three bash-command
instances per response; no single bullet produces more than one; removing
\texttt{solo-task} or \texttt{solo-explore} from \texttt{only-cr-imp}
drops out of the mode entirely. The multi-bullet response mode is not
cumulative pressure but a categorical phase transition in response
shape.

\begin{table}[H]
\centering
\caption{Bash-command-with-flag count per response on
\texttt{use-task-for-search} (Haiku, mean over 3 trials).
The emergent mode requires multiple bullets together.}
\label{tab:bash}
\small
\begin{tabular}{lc}
\toprule
Condition & bash-cmd-with-flag count \\
\midrule
\texttt{all-decl} (baseline) & 1.00 \\
\texttt{only-cr-imp} (all 7 bullets) & \textbf{3.00} \\
\midrule
\texttt{solo-explore} & 1.00 \\
\texttt{solo-task, solo-push, solo-dash-i, solo-no-edit,} & \\
\texttt{solo-empty, solo-heredoc, solo-empty-cr} & 0.00 \\
\bottomrule
\end{tabular}
\end{table}

% ============================================================
\section{Cross-Model Results}
\label{sec:crossmodel}

The E-SOLO protocol was replicated on Gemini 2.0 Flash and Claude
Sonnet 4.6. Neither reproduces the Haiku three-mode structure.

\subsection{Gemini: Ignores The Prohibition}

On every E-SOLO condition, Gemini invokes the Task tool on the
explore-agent probe. On \texttt{solo-task} --- where the CR block
explicitly contains ``NEVER use the TodoWrite or Task tools'' ---
Gemini opens ``I'll use the \texttt{`Task`} tool with the
\texttt{`Explore`} subagent\ldots'' and proceeds to emit
\texttt{Task(description=\ldots)}. The prohibition has no observable
effect on Gemini's response. Gemini's explore-agent score is 1.000
across all eight conditions.

This is not a replication of pathway A. It is a different failure
mode entirely: Gemini fails to be steered by the prohibition, whereas
Haiku is steered but into a different response mode.

On other probes Gemini exhibits its own idiosyncratic patterns
(e.g., producing \texttt{TodoWrite:} markdown code blocks on
\texttt{use-task-for-search} regardless of CR content). None of these
patterns match Haiku's register-bomb signatures.

\subsection{Sonnet: Default Sits At The Floor}

Sonnet's all-decl baseline on the explore-agent probe is 0.133 --- the
same level Haiku reaches under pathway A. Under every E-SOLO condition,
Sonnet stays at 0.10--0.20. There is no collapse because there is no
height to fall from. Sonnet's default response on this probe is a
prose-strategy response similar to Haiku's pathway-A output.

The one apparent outlier --- \texttt{solo-empty-cr} scoring 0.733 ---
is judge noise. Three trials scored 1.00, 1.00, and 0.20 on
essentially identical prose responses; the $0.20$ response differs
from the $1.00$ responses only in minor rephrasings that should not
justify a score gap. We attribute this to judge-level variability
on Sonnet, which we observed to be higher than on Haiku on this probe.

Sonnet does show real variation on \texttt{use-task-for-search}
(scores 0.20--0.60) accompanied by bash-command-with-flag counts
varying 1.0--3.3. Crucially, Sonnet's \emph{baseline} bash count is
3.33 --- higher than Haiku's super-additive peak of 3.00 in
\texttt{only-cr-imp}. Sonnet's CR bullets \emph{reduce} bash emission
from baseline; Haiku's (in multi-bullet interaction) \emph{increase}
it. Opposite pressures on the same probe.

\subsection{Three Model Patterns, No Shared Mechanism}

\begin{table}[H]
\centering
\caption{Cross-model summary of E-SOLO responses to register-isolated
CR content.}
\label{tab:crossmodel}
\small
\begin{tabularx}{\textwidth}{lXXX}
\toprule
Property & Haiku 4.5 & Gemini 2.0 Flash & Sonnet 4.6 \\
\midrule
Default explore-agent mode & Task-invoke (score 1.0) &
Task-invoke (score 1.0) & prose-strategy (score 0.13) \\
Response to ``NEVER use Task'' & Shifts to prose-strategy
(pathway A) & Ignores; invokes Task anyway & No visible effect
(already prose) \\
Response to empty CR block & Shifts to AskUserQuestion
(pathway B) & No effect & No effect \\
Multi-bullet interaction on\newline\texttt{use-task-for-search} &
Super-additive bash promotion & No bash-command emission &
Baseline-reduces bash emission \\
\bottomrule
\end{tabularx}
\end{table}

These three patterns do not compose into a single phenomenon.
The register-bomb collapse described in the prior paper
\citep{mason2026registerbombs} is a \emph{Haiku} phenomenon.
Gemini and Sonnet exhibit different phenomena that share only the
input (register-isolated CR content) but not the response shape or
mechanism.

% ============================================================
\section{Discussion}

\subsection{What The Prior Paper Measured}

The prior paper correctly identified that certain register-isolated
configurations catastrophically disrupt Haiku's adherence on
semantically unrelated probes. The experimental data in that paper
is unchanged. What E-SOLO reveals is that the \emph{mechanism
underlying the disruption} is not a single ``unbound prohibition''
but at least three response-mode phenomena (pathways A, B, and the
multi-bullet emergent mode), only one of which corresponds to what
the prior paper proposed. The ``register bomb'' is better read as
a behavioral family in Haiku than as a semantic mechanism.

\subsection{Why The Engineering Advice Still Works}

The prior paper's practical recommendations --- embed scope
per-prohibition, maintain register uniformity --- remain correct
for Haiku. Both interventions happen to avoid all three of Haiku's
pathways: inline-scoped prohibitions avoid pathway A (the clause
never achieves unbound reach); uniform imperative register across
blocks avoids pathway B (no register-contrast point exists to
destabilize response shape); using fewer imperative bullets per
block avoids the multi-bullet interaction mode.

The recommendations work for the right reasons only for pathway A.
They work for the wrong reasons (coincidentally) for pathways B and
the multi-bullet mode. This matters for anyone designing a prompt
engineering framework (or a DSL) aimed at eliminating these failure
modes: the naive approach of solving pathway A does not generalize.

\subsection{Scope Of The Mechanism Claim}

We urge that register-bomb claims be model-scoped in future work.
E-SOLO demonstrates that the phenomenon's existence is not an LLM
property but a Haiku-specific behavioral pattern. The
``probe transfer problem'' flagged in the prior paper's \S6 is
sharpened by this result: the probes did not transfer because the
phenomenon being probed is itself model-specific.

\subsection{E-SOLO As A Methodological Standard}

E-SOLO's design --- per-clause isolation paired with an empty-block
structural control --- is the minimum protocol for distinguishing
pathway A from pathway B. Prior single-experiment designs (including
those in the original paper) cannot separate them: the observed
score collapse in \texttt{only-cr-imp} conflates both. We recommend
that future register-bomb or unbound-prohibition claims be supported
by a solo condition for each clause under investigation plus an
empty-block null.

% ============================================================
\section{Implications For Prompt Governance}

The session-26 paper-correction note proposed renaming the
phenomenon to ``unbound prohibition'' on the grounds that this
names the mechanism more precisely than ``register bomb.'' We
modify that recommendation: ``unbound prohibition'' is the right
name for pathway A on Haiku. Pathway B requires a different name
(we suggest \emph{structural-ambiguity drift}) because no
prohibition is involved. The multi-bullet emergent mode requires
yet another name; we leave that open.

For a prompt governance framework
(e.g.~Arbiter~\citep{mason2026arbiter}) aiming to prevent these
failure modes at compile time, the following invariants are
suggested. These are necessary for Haiku; other models may need
different invariants.

\begin{enumerate}[leftmargin=*]
\item \textbf{Scope-welding (pathway A):} Prohibitions must weld
scope inline (``When X, NEVER Y''), not declare it in a prefix
(``During X: NEVER Y''). This is the invariant the prior paper
proposed.

\item \textbf{Content-sufficiency (pathway B):} Imperative-register
blocks must contain substantive content. An imperative heading
followed by no bullets --- or by bullets that do not constrain
behavior for most contexts --- can trigger structural-ambiguity
drift regardless of what the bullets say.

\item \textbf{Multi-clause interaction detection (emergent mode):}
Prohibitions whose named subjects have semantically adjacent
neighborhoods (e.g., Task and bash-for-code-reading) may jointly
trigger a response mode that no single clause produces. The
compiler cannot evaluate each clause in isolation; it must
also detect such co-occurrence.
\end{enumerate}

Of these, scope-welding and content-sufficiency are
machine-checkable. Multi-clause interaction detection is harder
and may require a behavior-model of the target LLM.

% ============================================================
\section{Limits}

\begin{enumerate}[leftmargin=*]
\item E-SOLO was run at three trials per condition. Mode assignment
is deterministic at $T=0$ within trial count; rare mode switches
or tail behaviors are not characterized.

\item The probe battery is text-only. Some response modes
(e.g.~Haiku's Task-invoke mode) involve markup that is never
executed. Whether the phenomenon replicates when tools can actually
be invoked is untested.

\item Cross-model comparison extends to three models (Haiku, Gemini,
Sonnet). DeepSeek, Mistral, and GPT-family models are not characterized.

\item The ``Important notes:'' heading structure used in pathway B
is specific to the Claude Code corpus. Whether structural-ambiguity
drift requires this specific structural signature or generalizes
to any imperative heading with insufficient content is untested.
\end{enumerate}

% ============================================================
\section*{Availability}

All experimental data, scripts, and analysis code are available at
\url{https://github.com/fsgeek/arbiter} under
\texttt{data/ablation/e\_solo/} and \texttt{scripts/run\_e\_solo.py}.
Reanalysis cairns T21, T22, T23 under \texttt{docs/cairn/}.

% ============================================================
\bibliography{references}

\end{document}
